% ChargeShield-FL --- DSN 2027 --- Preliminaries
% Sezione aggiunta 2026-09-16 (Sprint 10zz+113) su richiesta esplicita
% dell'utente: "quando spieghiamo il DP dovremmo dare anche le definizioni
% (idem per gli attacchi)" + "questa cosa dovremmo farla anche per i tre
% tipi di posizionamento del DP" + "dovremmo spiegare bene anche gli
% esperimenti fatti con no-DP".

\section{Preliminaries}
\label{sec:prelim}

\subsection{Differential Privacy}
\label{sec:prelim-dp}

Differential privacy~\cite{dwork2006calibrating,dwork2006ourdata,dwork2014algorithmic}
offers a guarantee for computations over aggregate datasets built on the
principle that, observing the output of a computation, one should not be
able to infer the presence or absence of any particular record in the
input. Formally, for two adjacent datasets $D$ and $D'$ differing in a
single record, the outcome of a differentially private mechanism on $D$
and on $D'$ should be near-indistinguishable.

\begin{definition}[$(\varepsilon,\delta)$-differential privacy]
\label{def:dp}
A randomised computation $\mathcal{A}$ satisfies
$(\varepsilon, \delta)$-differential privacy if, for any two adjacent
datasets $D$ and $D'$ and any subset $O \subseteq \mathrm{Range}(\mathcal{A})$,
\begin{equation}
\Pr[\mathcal{A}(D) \in O] \;\le\; e^{\varepsilon}\,\Pr[\mathcal{A}(D') \in O] \;+\; \delta .
\label{eq:dp}
\end{equation}
\end{definition}

The parameter $\varepsilon$ controls the trade-off between the accuracy of
the mechanism and the amount of information it releases, and is commonly
called its \emph{privacy budget}: a lower $\varepsilon$ means less leakage
and a stronger guarantee. The additive term $\delta$, introduced by
Dwork et al.~\cite{dwork2006ourdata}, admits the possibility that pure
differential privacy is violated with probability $\delta$, conventionally
chosen well below $1/|D|$. Our configuration uses $\delta = 10^{-5}$
against a client training set of $|D| \approx 1{,}344$ sessions, so
$\delta \ll 1/|D| \approx 7.4 \times 10^{-4}$ holds with margin.

\paragraph{What a record is here.} Definition~\ref{def:dp} is stated over
an adjacency relation, and the guarantee means nothing until that relation
is fixed. In our system clipping and noise are applied \textbf{once per
round to a client's aggregated update}, so the adjacency our mechanism
implements is \textbf{client-level}: $D$ and $D'$ differ by one
participating client. The attacks in Section~\ref{sec:prelim-mia}, by
contrast, probe \textbf{record-level} membership --- whether one charging
session was in training. These are different privacy units, and a
client-level guarantee does not by itself bound record-level leakage when
each client holds thousands of records, as it does here. We therefore
never read our $\varepsilon$ as a record-level bound on the quantity our
attacks measure; the relationship between the two is the empirical
question this paper studies, not an assumption it makes.

\subsection{Composition and Privacy Accounting}
\label{sec:prelim-accounting}

A single $(\varepsilon,\delta)$ guarantee does not describe a training run
of $T$ federated rounds, because each round releases a further
observation. Tracking the cumulative privacy loss across a sequence of
such releases is the task McSherry introduced as \emph{privacy
accounting}~\cite{mcsherry2009pinq}, and a composition theorem is what
bounds how the guarantee degrades under it.

We report three cumulative quantities per run. \emph{Basic composition}
gives $\varepsilon_{\text{tot}} = T \varepsilon$ and
$\delta_{\text{tot}} = T \delta$. \emph{Advanced
composition}~\cite[Thm.~3.20]{dwork2014algorithmic} gives, for any
$\delta' > 0$, a $(\varepsilon', T\delta + \delta')$ guarantee with
\begin{equation}
\varepsilon' = \varepsilon \sqrt{2T \ln(1/\delta')} \;+\; T\,\varepsilon\,(e^{\varepsilon} - 1),
\label{eq:advcomp}
\end{equation}
and we set $\delta' = \delta$. We then report
$\varepsilon_{\text{best}} = \min(\varepsilon_{\text{tot}}, \varepsilon')$.

Two points deserve emphasis, because both are easy to misread. First,
advanced composition is \textbf{not uniformly tighter}: its leading term
grows as $\sqrt{T}$ rather than $T$, which is an improvement only for
large $T$ and small $\varepsilon$. At our operating point
($\varepsilon = 1.0$, $\delta = 10^{-5}$, $T = 3$)
Equation~\ref{eq:advcomp} yields $\varepsilon' \approx 13.47$ against
$\varepsilon_{\text{tot}} = 3.0$, so the basic bound is the one that
holds and $\varepsilon_{\text{best}} = 3.0$. We report the looser number
alongside it rather than silently dropping it.

Second, and more consequential: we perform composition accounting, but we
do \textbf{not} implement a privacy accountant in the
moments-accountant/R\'enyi sense~\cite{abadi2016deep}, and we do not
\emph{enforce} a budget. Nothing in our system halts training when a
cumulative budget is reached --- unlike PINQ~\cite{mcsherry2009pinq},
which refuses queries once the budget is exhausted. Our Privacy Auditor
(Section~\ref{sec:framework-auditor}) maintains a per-client running
$\varepsilon$ and raises alerts, but that quantity is a sensitivity-scaled
operational heuristic for anomaly detection, not formal accounting, and it
is deliberately not the figure we report.

The deepest caveat is about the per-round quantity being composed. We clip
and add noise once per round to a client-level update, with $50$ local
epochs inside the round and no per-sample gradient clipping. This is not
DP-SGD, so the per-round $(\varepsilon,\delta)$ is a configuration label
calibrating the noise scale rather than a verified per-round guarantee.
Composing a label yields a composed label. We state this explicitly
because the alternative --- reporting $\varepsilon_{\text{best}} = 3.0$
without it --- would claim a formal guarantee our mechanism does not
establish.

\subsection{Three DP Placements as Observation Points}
\label{sec:prelim-placements}

Let $g_c^{(t)}$ be the raw update produced by client $c$ in round $t$.
The privatisation pipeline applies, in order, clipping to $L_2$ norm $C$
and additive Gaussian noise:
\begin{equation}
g \;\xrightarrow{\;\text{clip}\;}\; \bar{g} = g \cdot \min\!\left(1, \tfrac{C}{\lVert g \rVert_2}\right)
\;\xrightarrow{\;\text{noise}\;}\; \tilde{g} = \bar{g} + \mathcal{N}(0, \sigma^2 I).
\label{eq:pipeline}
\end{equation}

Our three configurations --- \texttt{dp-fedavg}, \texttt{central},
\texttt{local} --- are \textbf{not three privacy mechanisms}. The
mechanism is the same Gaussian mechanism in all three. What differs is
\emph{where along Equation~\ref{eq:pipeline} the adversary is assumed to
observe}: \texttt{dp-fedavg} observes $g$, \texttt{central} observes
$\bar{g}$, and \texttt{local} observes $\tilde{g}$. They are three points
of adversary privilege on one pipeline, ordered by how much
transformation has already been applied.

This ordering is what makes the design informative.
\texttt{dp-fedavg} is the \textbf{upper bound} on what any adversary can
extract from this channel, because it sees the update before any
protective transformation. If the most privileged view yields no
membership signal, then a null result under the two protective placements
cannot be attributed to noise having suppressed a signal that was
otherwise present: there was none to suppress. Reporting the three as if
they were competing mechanisms would lose exactly this argument.

\subsection{Membership Inference Attacks}
\label{sec:prelim-mia}

A membership inference attack asks, for a target record $x$ and a model
$f$ trained on an unknown dataset $D$, whether $x \in D$. We follow the
standard formulation~\cite{shokri2017membership,yeom2018privacy}: the
adversary computes a score $s(x, f)$ intended to be stochastically larger
for members than for non-members, and membership is decided by
thresholding $s$. We report AUC-ROC, which is threshold-free and equals
$\Pr[s(x_{\text{mem}}) > s(x_{\text{non}})]$ over independently drawn
member/non-member pairs; $0.5$ is chance.

Our target is an \textbf{unsupervised} autoencoder, so the per-record
signal is the reconstruction loss
$\ell(x, f) = \lVert f(x) - x \rVert_2^2 / d$ rather than a classification
loss. We instantiate three scores:

\begin{itemize}
\item \textbf{Yeom}~\cite{yeom2018privacy}: $s = -\ell(x,f)$. A member is
predicted when the loss is low. It requires no auxiliary training and
uses no calibration.

\item \textbf{Shadow}~\cite{shokri2017membership,salem2019mlleaks}: a
shadow model $f_s$ is trained on a disjoint split, and its loss
distribution supplies the threshold applied to $\ell(x,f)$. This
calibrates \emph{globally} --- one reference distribution for all records.

\item \textbf{LiRA}~\cite{carlini2022membership}: calibrates
\emph{per-record}. For each $x$, $N$ shadow models are partitioned into
those trained with $x$ (IN) and without it (OUT), giving Gaussian fits
$\mathcal{N}(\mu_{\text{in}}, \sigma_{\text{in}}^2)$ and
$\mathcal{N}(\mu_{\text{out}}, \sigma_{\text{out}}^2)$; the score is the
log-likelihood ratio
$s = \log p(\ell \mid \text{IN}) - \log p(\ell \mid \text{OUT})$.
\end{itemize}

Yeom and Shadow observe the aggregated global model, so a positive result
places a record in \emph{some} client's training data. LiRA observes the
per-client update, so a positive result attributes membership to a
specific operator, but requires aggregator privilege. These are two
independent observation surfaces rather than a hierarchy, and we treat
them as such throughout.

\subsection{The no-DP Arm}
\label{sec:prelim-nodp}

Every configuration is also run with the noise disabled ($\sigma = 0$),
and this arm does more work than a baseline label suggests. A near-chance
AUC under DP is ambiguous on its own: it is consistent both with DP having
suppressed a real signal and with no extractable signal having been
present. The no-DP arm resolves that ambiguity in one direction. If the
attack is at chance \emph{without} any noise, on the raw update, then a
null result with noise is not evidence that DP worked --- there was
nothing to protect against in the first place, and any claim of protective
effect would be unfalsifiable. Conversely, a measurable no-DP signal is
what makes a subsequent drop under DP interpretable as mitigation.

The no-DP arm is therefore a precondition for reading the DP arms at all,
not a control we report for completeness, and we present it first in
Section~\ref{sec:results} for that reason.
